% ===========================================================================
\section{Why it always exists}\label{sec:existence}
\warmth{0}\quad\whisper{Eigen-decomposition can fail; the SVD never does. The reason is $A\T A$.}

\begin{strand}
Eigenvalues are fussy: a matrix may not be diagonalizable, may have complex eigenvalues, must
be square. The SVD has \emph{none} of these defects. The trick is to launder $A$ through the
one matrix that is always perfectly behaved: $A\T A$ is symmetric and positive semidefinite, so
the \keyword{spectral theorem} hands us an orthonormal eigenbasis on a silver platter. The
right singular vectors are those eigenvectors; the singular values are square roots of the
eigenvalues.
\end{strand}

\block{Existence, in one construction}

\begin{theorem}[every matrix has an SVD]\label{thm:exist}
Every $A\in\R^{m\times n}$ admits a singular value decomposition $A=U\Sig V\T$.
\end{theorem}
\begin{proof}
Skeleton (carry it out as you read):
\begin{enumerate}[leftmargin=2em,itemsep=1pt]
  \item $A\T A$ is symmetric PSD, so by the spectral theorem $A\T A=V\Lambda V\T$ with $V$
        orthonormal and $\Lambda=\diag(\lambda_1\ge\cdots\ge\lambda_n\ge 0)$.
  \item Set $\sigma_i=\sqrt{\lambda_i}$. For each $i$ with $\sigma_i>0$, define
        $u_i=\dfrac{1}{\sigma_i}\,A v_i$. \TODO{check $\{u_i\}$ is orthonormal, using $v_i\T A\T A v_j=\lambda_j\,v_i\T v_j$}
  \item Extend $\{u_i\}$ to an orthonormal basis $U$ of $\R^m$. \TODO{verify $A=\sum_i\sigma_i u_i v_i\T = U\Sig V\T$ by checking both sides agree on the basis $\{v_j\}$}
\end{enumerate}
\end{proof}

\recall{Mirror route: $AA\T=U\Sig\Sig\T U\T$ gives the \emph{left} vectors $u_i$ directly. $A\T A$ and $AA\T$ share the nonzero eigenvalues $\sigma_i^2$.}

\begin{remark}[how unique?]
The singular \emph{values} $\sigma_i$ are unique. The singular \emph{vectors} are not: each
$(u_i,v_i)$ may be flipped in sign together, and within a repeated singular value you may rotate
the corresponding vectors. So $\Sig$ is rigid; $U,V$ have a little freedom.
\end{remark}

\begin{yourturn}
Compute the SVD of $\displaystyle A=\begin{pmatrix}0&2\\3&0\end{pmatrix}$ by the recipe above.
\[
  A\T A=\fillin[2.5cm],\quad \text{eigenvalues } \fillin[1.6cm]\Rightarrow \sigma_1,\sigma_2=\fillin[1.4cm],
\]
\[
  v_1=\fillin[1cm],\ v_2=\fillin[1cm],\qquad
  u_i=\tfrac{1}{\sigma_i}Av_i:\ u_1=\fillin[1cm],\ u_2=\fillin[1cm].
\]
\recall{$A\T A=\diag(9,4)$, so $\sigma=3,2$, $v_i=e_i$, and $u_1=e_2,u_2=e_1$. Notice $U\ne V$: this map is a reflection-and-stretch, not a symmetric one.}
\end{yourturn}

\begin{strand}
When $A$ \emph{is} symmetric PSD, the SVD and the eigendecomposition coincide ($U=V$,
$\sigma_i=\lambda_i$). For symmetric but indefinite $A$, $\sigma_i=|\lambda_i|$ and a sign moves
into $U$. The SVD is the eigendecomposition's well-behaved cousin.
\end{strand}
